% QR-Code für LP-06-reibung-vertiefung
\documentclass[a4paper]{article}
\usepackage[margin=2cm]{geometry}
\usepackage{qrcode}
\usepackage[hidelinks]{hyperref}
\usepackage{xcolor}
\renewcommand{\familydefault}{\sfdefault}

\definecolor{physik}{HTML}{2c5aa0}

\begin{document}
\pagestyle{empty}

\begin{center}
\vspace*{3cm}

{\LARGE\bfseries\textcolor{physik}{Lernpfad Physik Klasse 7}}

\vspace{1cm}

{\Huge\bfseries Reibung Vertiefung}

\vspace{0.5cm}

{\large ABS \& Anwendungen}

\vspace{2cm}

\href{https://devbox42.github.io/sciencesim/physik/kl07/02-kraefte/06-reibung-vertiefung/LP-06-reibung-vertiefung.html}{%
\qrcode[height=8cm]{https://devbox42.github.io/sciencesim/physik/kl07/02-kraefte/06-reibung-vertiefung/LP-06-reibung-vertiefung.html}%
}

\vspace{1.5cm}

{\small\url{https://devbox42.github.io/sciencesim/physik/kl07/02-kraefte/06-reibung-vertiefung/LP-06-reibung-vertiefung.html}}

\end{center}
\end{document}
